\documentclass[10pt,a4paper]{amsart}
\usepackage[margin=19mm]{geometry}
\usepackage{amsmath,amssymb,mathtools}
\usepackage[scr=rsfs,cal=euler]{mathalfa}
\usepackage{xfrac}
\usepackage[unicode,hidelinks,bookmarks=false]{hyperref}
\newcommand{\ie}{i.e.\ }
\newcommand{\cf}{cf.\ }
\newcommand{\resp}{respectively\ }
\newcommand{\Subsection}{\subsection}
\providecommand{\nobreakdash}{\nobreak\hbox{-}}
\setlength{\emergencystretch}{2em}
\multlinegap0pt
\usepackage{bm}
\usepackage{bbm}
\usepackage{dsfont}
\usepackage{xspace}
\usepackage{cancel}
\usepackage{mathtools}
\usepackage{amsthm}
\makeatletter
\@ifundefined{theorem}{\newtheorem{theorem}{Theorem}}{}
\@ifundefined{claim}{\newtheorem{claim}[theorem]{\bf Claim}}{}
\makeatother
\makeatletter
\expandafter\ifx\csname proof\endcsname\relax
\renewenvironment{proof}[1][\proofname] {\par\pushQED{\qed}\normalfont\topsep6\p@\@plus6\p@\relax\trivlist\item[\hskip\labelsep\textit{#1}\@addpunct{.}]\ignorespaces}{\popQED\endtrivlist\@endpefalse}
\fi
\makeatother
\title[Cyclic subsets of tournaments]{Cyclic subsets of tournaments}
\author{Zach Hunter, Teng Liu, Aleksa Milojevi\'c, Benny Sudakov}
\date{Source: arXiv:2508.03634v1 (CC BY 4.0)}
\begin{document}
\maketitle

\noindent\emph{Source and licence.} Cyclic subsets of tournaments. Version arXiv:2508.03634v1 (CC BY 4.0), distributed under the Creative Commons Attribution 4.0 International licence, as recorded in the arXiv licence metadata for this exact version and shown on the abstract page https://arxiv.org/abs/2508.03634v1. Full rights notice: licenses/arxiv-2508.03634-CC-BY-4.0.txt.\par\smallskip

\noindent\textit{Editorial scope.} Counted unit: the Claim whose source label is \texttt{claim:stability} in the article (source0.tex lines 391--393, source label \texttt{claim:stability}), delivered together with the article's own complete original proof of it (source0.tex lines 394--409, one \texttt{proof} environment of 16 lines). No mathematics is written, repaired or re-derived: statement and proof are the article's own text line for line. Required context retained uncounted: none. Every definition, display and label of those lines is retained unchanged. Omitted from this package and not redistributed: 1--390, 410--670. Nothing inside a retained excerpt is omitted or altered: every retained body is delivered exactly as the article's lines are written, so no omission receipt of any kind accompanies this package. Citations: the retained excerpts carry no citation command at all, so no bibliography entry of the article is transcribed and no reference is redistributed. Reference adaptation: none was needed and none was made; every reference inside the delivered unit resolves inside it, so no unresolved reference or citation is produced. Printed slips kept literal: no printed word, symbol, hypothesis or number of the counted statement or of its proof was altered. Layout and class: the article's own class is \texttt{article} with packages including amsmath, amssymb, amsthm, graphicx, amscd, epic, eepic, url, xcolor, inputenc, comment, enumerate, epstopdf, enumitem; the wrapper is the lane's pinned \texttt{amsart} editorial wrapper at 10pt on A4, which restates the theorem-like environments the retained text needs and adds the article's own macro definitions that the retained text needs (none), with extra wrapper packages (bm, bbm, dsfont, xspace, cancel, mathtools). The delivered numbering is the unit's own. Assets: no retained span contains a figure environment, a graphics-inclusion command, a PDF-page inclusion, a table or a TikZ picture; no image file is copied into this package, the package declares no asset dependency and no image file is redistributed with it. Licence: version arXiv:2508.03634v1 of this article is distributed under the Creative Commons Attribution 4.0 International licence, as recorded in the arXiv licence metadata for this exact version and shown on the abstract page https://arxiv.org/abs/2508.03634v1. A full rights notice is delivered at licenses/arxiv-2508.03634-CC-BY-4.0.txt. Characters: every retained excerpt is pure ASCII, so the delivered engine, XeLaTeX with its default Latin Modern text font, reports no missing character.
\begin{claim}\label{claim:stability}
Let $\delta\in (0, 1/2)$ be a constant and let $T$ be an $n$-vertex tournament with $\delta^0(T)\geq (1/4-\delta^2)n$ with no Hamilton cycle. Then, $T$ contains at least $(1/2-2\delta) n$ vertices of in-degree smaller than $(1/4+2\delta)n$. 
\end{claim}

\begin{proof}
If $T$ is not Hamiltonian, it is not strongly connected and so there exists a directed cut $(A, B)$. Since each vertex $v\in V(T)$ satisfies $d^+(v)\geq n/4-\delta^2 n$ and there are no edges going from $B$ to $A$, we have that 
\[\binom{|B|}{2}=\sum_{v\in B}d^+(v)\geq |B|\cdot (n/4-\delta^2 n),\]
which shows that $|B|\geq n/2-2\delta^2 n$. Thus, $|A|=|V(T)|-|B|\leq n-n/2+2\delta^2 n=n/2+2\delta^2 n$. 

Similarly, by computing the sum of in-degrees of the vertices in $A$ and using that $d^-(v)\geq n/4-\delta^2 n$ for all $v\in A$, we find that $\binom{|A|}{2}=\sum_{v\in A}d^-(A)\geq |A|\cdot (n/4-\delta^2 n)$, which shows that $|A|\geq n/2-2\delta^2 n\geq n/2-\delta n$ (since $\delta\leq 1/2$).% Thus, $|A|=|V(T)|-|B|\leq n-n/2+2\delta^2 n=n/2+2\delta^2 n$.

Let $r$ be the number of vertices $v\in A$ with $d^-(v)\geq (1/4+2\delta) n$. To finish the proof, we need to show that $r\leq \delta n$, since this implies that at least $|A|-\delta n\geq n/2-2\delta n$ vertices have $d^-(v)\leq (1/4+2\delta) n$. 

Let us consider again the sum of in-degrees of vertices $v\in A$, now taking into account the $r$ vertices of high in-degree. %On one hand, we have $\sum_{v\in A}d^-(v)=\binom{|A|}{2}$, since every edge is counted exactly once and no edge goes from $B$ to $A$. 
Since every vertex of $A$ has in-degree at least $\delta^0(T)\geq (1/4-\delta^2)n$, and there exist $r$ vertices which have additional $2\delta n$ incoming edges, we have
\[\binom{|A|}{2} = \sum_{v\in A}d^-(v)\geq |A| (n/4-\delta^2 n) + r\cdot 2\delta n.\]
Rearranging, we find that 
\[r\leq \frac{1}{2\delta n}\cdot |A|\left(\frac{|A|-1}{2}-\frac{n}{4}+\delta^2 n\right) \leq \frac{1}{2\delta n}\cdot (1/2+2\delta^2) n \cdot 2\delta^2 n  \leq \delta n,\]
where we have used that $|A|\leq n/2+2\delta^2 n$ in the second inequality. This shows $r\leq \delta n$ and thus completes the proof.
\end{proof}

\end{document}
